\subsubsection{Convex Hull Trick}

\prob{1}{dp-cvxhull-trick1} Suponha que você tenha um conjunto de funções lineares, isto é, $S = \{kx + c | k, c \in \mathbb{Z}\}$. Você pode fazer duas coisas: adicionar funções ao conjunto e fazer queries que pedem, dado $x$, a função $f$ que minimiza $f(x)$.

\textbf{Solução:} Interprete as funções como pontos $(k, c)$ no plano. Dado uma query $x$, você quer encontrar o ponto $(k, c)$ que minimiza $(x, 1) \cdot (k, c)$. É possível ver que os pontos que podem ser solução formam a parte de baixo de um fecho convexo. A intuição da resposta é pegar a reta normal a $(x, 1)$ (que tem inclinação $-x$), e ver o primeiro ponto que ela toca. As arestas do fecho estão ordenadas em relação à sua inclinação, então basta encontrar o ponto cuja inclinação das arestas passa $-x$ (isto é, a aresta anterior do ponto é menor que $-x$, e a próxima é $\ge -x$). Adicionar novos pontos deve ser feita usando uma estrutura de dados ordenada que permite inserção e remoção arbitrária (tipo set) para que as adições possam ser feitas em qualquer ordem.

\inccode{dp/convex_hull_trick.cpp}

\prob{2}{dp-cvxhull-trick2} Uma versão mais simples do problema anterior, mas as adições são com $k$ crescente e as consultas podem ser feitas num intervalo da convex hull.

\textbf{Solução:} Solução parecida com a anterior, mas dessa vez se utiliza um vetor para no lugar do set para permitir as queries em intervalo. Como as adições são ordenadas, apenas o final do vetor é modificado.

\inccode{dp/convex_hull_trick_simples.cpp}

